\chapter{Conclusion}\label{ch:conclusion}

\section{What worked}\label{sec:conclusion-what-worked}

SR-FCA did exactly what the paper promised. On both MNIST-rotated ($K^\star = 4$) and MNIST-inverted ($K^\star = 2$) it recovered the ground-truth clustering on every seed, with zero misclustering and the correct $K^\star$ discovered without being told. The per-cluster test accuracies -- $93.02 \pm 0.17\,\%$ on rotated and $92.95 \pm 0.18\,\%$ on inverted -- land slightly above the paper's reported $91.66 \pm 0.13$ and $92.03 \pm 0.30$ respectively. The same qualitative behaviour transferred to Fashion-MNIST in both rotation-based ($K^\star = 4$) and inversion-based ($K^\star = 2$) splits, confirming that the clustering primitive is robust to the source of heterogeneity.

The pairwise-distance-histogram procedure for choosing $\lambda$ was the most pleasant surprise of the reproduction. It is cheap (one round of local training), it produces an interpretable artifact (the histogram), and it generalises across datasets: the same valley-midpoint heuristic produced a working $\lambda$ on every (dataset, split) pair we tried, including the K-pair Fashion-MNIST splits the paper did not study. Compared to the effort of picking $K$ for IFCA -- running the whole pipeline multiple times -- this is a concrete ergonomic win.

\section{What did not work well}\label{sec:conclusion-did-not-work}

The clearest empirical limitation we observed is that on clean simulated data, ONE\_SHOT alone already recovers the correct clustering -- so RECLUSTER and MERGE are no-ops in subsequent REFINE rounds, and the trimmed-mean aggregator's robustness margin is never exercised. The A1 ablation makes this explicit: the cluster \emph{assignment} is perfect at $R = 0$, and the $22.25 \to 91.72 \to 93.08\,\%$ accuracy progression across $R = 0/1/2$ comes entirely from training the per-cluster models, not from correcting the clustering. Theoretically the robustness guarantee (Theorem~4.14) is what makes SR-FCA's correctness claim go through; empirically, on idealised image-rotation/inversion benchmarks, plain FedAvg inside each cluster would have produced the same numbers. The value of trimming should show up on real heterogeneous data with Byzantine or straggler clients, but that experiment is outside our scope.

A second mild limitation is that we did not re-train the five baseline methods (Local, Global, CFL, LOCAL-KMeans, IFCA) at the paper's full configuration; we quote their numbers from Tables~2 and~3 of~\parencite{vardhan2024srfca} instead. Re-implementing each at $T = 280$ would have multiplied the compute budget by~$6$. The cross-method ordering in our report therefore inherits the paper's, rather than being independently verified.

\section{Open questions}\label{sec:open-questions}

Three follow-ups we would like to see, already flagged in \S\ref{sec:future}. First, a hybrid SR-FCA$\to$IFCA that uses SR-FCA for cluster discovery and then hands off to IFCA for the faster $\tilde{\mathcal{O}}(1/n)$ per-cluster rate. Second, a scalable ONE\_SHOT via Johnson--Lindenstrauss sketching, to address the $\mathcal{O}(m^2)$ pairwise cost at $m \gg 10^3$ scale. Third, a comparison of SR-FCA against FedSoft and FedProx aggregation inside REFINE on \emph{real} federated data -- we covered IFCA and Local-KMeans in this report because they are the closest competitors to SR-FCA per the paper, but an end-to-end comparison with FedSoft was outside scope.

\section{What we learned}\label{sec:lessons}

Three lessons for future reproductions of clustered-FL methods. (i) The pairwise-distance histogram is underrated -- it is cheap, it answers the hardest hyperparameter question ($\lambda$, or implicitly $K$), and it is a free diagnostic for whether the clustering assumption holds on a given dataset. (ii) Ablation designs should include the axis ``what part of the algorithm is doing the work?'' -- on clean simulated data the A1 ablation showed that REFINE's contribution is entirely in training the cluster models, with cluster recovery already complete after ONE\_SHOT, which is more useful than re-confirming an effect that is obvious from theory. (iii) When the paper's full experimental grid does not fit a course's compute budget, reproducing one method at the paper's actual configuration and quoting the rest of the comparison is more honest than reproducing every method at a reduced configuration: the absolute numbers in the former case are interpretable against the paper, whereas the latter can only support relative claims.
